\section{Movement}\label{sec:6.0}

\ruleslabel{General Rule}

\begin{generalrule}
After their deployment on the map, Axis and Yugoslav units may move either within the Zone they currently occupy or from one Zone to an adjacent Zone. At all times during the game, units on the map must occupy either the Zone Display or an Objective Display. Exception: See 5.4.
\end{generalrule}

\ruleslabel{Cases}

\subsection{How To Move A Unit}\label{sec:6.1}

During a Movement Phase, the Phasing Player may move his units in any order he wishes. Eligible units are moved individually; once a Player begins to move a given unit, he must complete its movement before any other unit is moved. During a Movement Phase, all, some, or none of a Player’s units may be moved.

\case{6.11} If a Player wishes to move a unit, he may move it within the Zone it currently occupies or to any adjacent Zone. An ‘‘adjacent’’ Zone is any Zone with a common border of any length with the Zone currently occupied by a unit.

\case{6.12} Within a given Zone, a unit must occupy the Zone Display or any of that Zone’s Objective Displays.

\case{6.13} If a unit is positioned in a Zone Display, it must be deployed into either the Mountains or the Hide-away. In addition, the unit must be positioned in a box analogous to its nationality (e.g., Yugoslav/Partisan).

\case{6.14} If a unit is positioned in an Objective Display, it must be deployed in the box analogous to its type (e.g., Axis).

\case{6.15} A Yugoslav unit may never be placed in an Axis box, circle, or triangle and vice versa. However, the presence of an Enemy unit in a corresponding location has no effect on a unit’s movement capability.

\case{6.16} When moving within a Zone, a unit may be placed in any of that Zone’s Objective Displays or the Zone Display. Exception: See 6.34.

\subsection{Stacking Restrictions}\label{sec:6.2}

\case{6.21} No box, circle, or triangle may ever contain more than seven divisions or their equivalent at the end of a Movement Phase (see 6.22).

\case{6.22} A battalion is considered one-sixth of a division, a regiment is considered two-sixths (i.e., one-third) of a division, and a brigade is considered three-sixths (i.e., one-half) of a division. Thus, if a battalion, regiment, and brigade occupied the same location, this would be considered the equivalent of one division. All Soviet corps are considered the equivalents to divisions.

\case{6.23} Units of different nationalities may stack in the same location as long as they are friendly to each other (see 3.0). Enemy units and Partisan and Chetnik units may never stack together.

\case{6.24} Guerrilla groups (see 7.6) are stacked in boxes, circles, or triangles for ‘‘free.’’ An unlimited number of these units may occupy the same location.

\subsection{Axis Movement Restrictions}\label{sec:6.3}

\case{6.31} Axis units are permitted to move three contiguous Zones per Axis Movement Phase. Exception: See 6.32. Example: An Axis unit in Baranya could move to Serbia or Bosnia. Assuming it moves to Bosnia, it could then enter Carinthia, Croatia, or Serbia. Assuming it enters Carinthia, it could complete its movement by entering Istria, Slovenia, or Croatia.

\case{6.32} Pro-Axis Chetnik units (see 7.2) are moved exactly like Yugoslav units (see 6.4).

\case{6.33} When the final Zone an Axis unit enters in a Movement Phase has been decided upon, the unit must be positioned in the Zone Display or any Objective Display of that Zone.

\case{6.34} Axis units may never enter a Hide-away of a Zone Display except when on Anti-Guerrilla Operations (see 8.4.)

\subsection{Yugoslav Movement Restrictions}\label{sec:6.4}

\case{6.41} All Yugoslav units (including Tito, Partisan, pro-Yugoslav Chetnik, Soviet, and pro-Yugoslav Bulgarian units) are permitted to move within the Zone they currently occupy or to any adjacent Zone during each Yugoslav Movement Phase. Example: A Partisan unit occupying Macedonia could move within Macedonia or to Montenegro or Serbia.

\case{6.42} When a Yugoslav unit enters a Zone during the Yugoslav Movement Phase, it must end this Movement Phase in the Mountains or the Hide-away of the Zone Display — it may not move directly into an Objective Display. Exceptions: See 6.43 and 10.36. Example: A Yugoslav unit moving from Macedonia into Serbia must end the Movement Phase in the Mountains or Hide-away of the Serbian Zone Display. It could not move directly into a Serbian Objective Display.

\case{6.43} Soviet and pro-Yugoslav Bulgarian units are not subject to the restrictions of Case 6.42. When such a unit enters a Zone, it may be positioned in the Zone Display or any of the Zone’s Objective Displays. However, note that Soviet and pro-Yugoslav Bulgarian units are subject to specific movement restrictions (see 6.65).

\case{6.44} Soviet and pro-Yugoslav Bulgarian units should be considered Partisan units for movement and combat purposes. When occupying a Zone Display, they must stack in the Yugoslav/Partisan triangles or circles. They may stack freely with Partisan units and may participate in combat with them without restriction. Soviet and pro-Yugoslav Bulgarian units, however, may never stack with Chetnik units.

\subsection{Islands}\label{sec:6.5}

A unit may enter the Islands Zone only from Istria, Slovenia, Croatia, or Dalmatia. The Islands Zone consists only of a single Zone Display — it does not possess any Objective Displays.

\subsection{Occupation Zone Movement Restrictions}\label{sec:6.6}

Units of different nationalities may be restricted in the Zones that they may enter or move through during the course of the game. If a unit is prohibited from entering a given Zone, it is also prohibited from retreating after combat into that Zone. For a summary of these restrictions, see the Occupation Zone Table (6.7).

\case{6.61} Between Game-Turns 1 and 14 (inclusive), Partisan and Chetnik units may enter any Zone on the map except Carinthia and Baranya. Exception: See 6.63. Starting with Game-Turn 15, they may enter any Zone on the map.

\case{6.62} Each Axis nationality has a ``sphere of influence'' limited to the Zones listed below. An Axis unit may never enter a Zone outside of its sphere of influence. Exception: See 6.64 and 6.66. However, beginning with Game-Turn 15, all movement restrictions are lifted for German units — they may enter any Zone on the map.

\textbf{A.} German: Serbia, Bosnia;

\textbf{B.} Italian: Albania, Croatia, Dalmatia, Islands, Istria, Montenegro, Slovenia;

\textbf{C.} Bulgarian: Macedonia (see also 6.65);

\textbf{D.} Croatian: Bosnia, Croatia (see also 6.66); E. Serbian: Serbia;

\textbf{F.} Pro-Axis Chetnik: See Case 6.61.

\case{6.63} On Game-Turns 1 and 2, all units of both Players must remain in the Zone in which they started the game or in which they entered the game as reinforcements.

\case{6.64} During all Axis Movement Phases following the Game-Turn in which the Yugoslav Player first accumulates 45 Victory Points (or in which the first Yugoslav brigade has been placed on the map), the following movement restrictions are lifted:

\textbf{A.} German units may enter all Zones in the Italian sphere of influence;

\textbf{B.} Up to three Italian divisions may enter Serbia or Bosnia on an Anti-Guerrilla Operation (see 8.4), but not if an Italian Pullback has occurred (see 10.2). However, if Italian units participate in an Anti-Guerrilla Operation in Serbia or Bosnia, they must redeploy into Croatia (where they are considered part of the Axis garrison; see 7.5).

\textbf{C.} The Bulgarian 25th and 27th Divisions may enter any Objective Display in Serbia labeled with a ``Bulgarian Occupation'' in the Axis box. In addition, these units may participate in an Anti-Guerrilla Operation in Serbia, but are obligated to redeploy after the Operation into a box labeled with a ``Bulgarian Occupation.''

\case{6.65} Soviet and pro-Yugoslav Bulgarian units may not move out of Serbia
or Macedonia on Game-Turn 15.

\case{6.66} In all Axis Movement Phases following Italian Surrender (see 10.3), Croat units may include Dalmatia as part of their sphere of influence.

\case{6.67} Between Game-Turns 1 and 13 (inclusive), there may be no more than four Partisan and four Chetnik units (of any allegiance) in Macedonia. The Players may not move Partisan and Chetnik units into this Zone in excess of these limitations, nor may guerrilla reinforcements (see 7.4) be created in Macedonia that would violate this restriction.

\subsection{Occupation Zone Table}\label{sec:6.7}

(see mapsheet)
